\section{Limitations and threats to validity}
\label{sec:limitations}

We state these plainly, because several of them bound how strongly the headline
comparison can be read.

\paragraph{The reference policy is not a competent model.}
GRPO's pinned $\polref$ is the policy at initialization: an untrained MLP
adaptor head with a $-10$ logit bias, i.e.\ a near-uniform, high-entropy prompt
distribution --- not a supervised-fine-tuned model as in canonical RLHF
pipelines. Pinning such an anchor is a genuine departure from canonical GRPO.
It is anti-collapse, so it works in GRPO's favour here, but it also means the KL
term pulls toward high entropy rather than toward a good policy, and a GRPO run
starting from a competent reference might behave differently. This is the single
most important caveat on the comparison.

\paragraph{The tradeoff curve itself is not yet measured.}
As Sections~\ref{sec:tradeoff} and~\ref{sec:provenance} explain, the per-$\lmb$
frontier --- the figure this paper is nominally about --- is pending a rerun.
Everything reported here is $\lmb$-averaged. It remains possible that two arms
with similar average reward have differently-shaped curves, and the dominance
claim of Figure~\ref{fig:tradeoff} is about averaged operating points, not about
pointwise dominance across $\lmb$.

\paragraph{The rerun is single-seed.}
The regeneration campaign trains one seed per arm, not three. Its frontiers will
therefore be indicative rather than seed-averaged, and given GRPO's wide seed
spread ($26.3$--$34.7$ at step $11{,}500$) a single GRPO seed should not be read
as that arm's typical curve.

\paragraph{Evaluation noise and the number of seeds.}
Three seeds per arm and a $1.0$-point noise floor are enough to separate the
families ($6.0$ points) but not to rank variants within a family; we
consequently report the top two \rrebel{} variants as tied and draw no
conclusion between them. The noise floor is itself estimated from a single
repeated evaluation, so it is an order-of-magnitude figure rather than a
calibrated standard error.

\paragraph{The style classifier used at evaluation.}
The evaluation entry point requests the \emph{train}-split style classifier
rather than the test-split one. Since the same classifier scores every arm, this
cannot explain the family gap, but it does mean the absolute sentiment numbers
are not directly comparable to work that evaluates with the held-out
classifier.
\todo{confirm whether this is intentional; if not, re-evaluate with the
test-split classifier and update the absolute sentiment numbers. Arm ordering is
not at risk either way.}

\paragraph{Single task and single direction.}
All results are Yelp negative-to-positive with a \texttt{gpt2-xl} task model and
$5$-token prompts. We make no claim about other datasets, other transfer
directions, longer prompts, or larger task models. In particular, the inert-clip
argument of Section~\ref{sec:grpo} is a property of the single-update regime; a
pipeline that takes multiple epochs per rollout would keep GRPO's trust region
live and might close some of the gap.

\paragraph{Scope of this draft.}
Per its stated scope, this draft presents the algorithms and the results.
Related work is a stub, and there is no discussion or conclusion section yet.
